\section{Expansión de la acción exterior}
Todas las estimaciones holomorfas se toman en discos fijos; se puede disminuir sus radios
para diferenciar los restos. Los radios en $x$ siguen siendo mayores que uno antes de
evaluar vértices. En los parámetros físicos,
\[
 R(x)={\rm af\acute in}(x)+s^{-1}\mathcal R(\sigma,h,x).
\]
En efecto, sustituir $p=s t_rt_c v$ en la entropía permite cancelar los términos
$\sum t_v\log t_v$ con los márgenes. Los términos en $\log s$ y en $\log z$
son afines en $x$ por factorización de $z$ en su soporte y por los márgenes exactos.
En la solución del apartado anterior se evalúa
\[
 \mathcal R=\int\!\left[-v\log(v/z)-v+
 \sum_{l\ge2}\frac{\sigma^{l-1}(t_rt_c)^{l-1}v^l}{l(l-1)}\right]dM_1^R dM_1^C.
\]
Se usa $\log q$ en el primer sumando y el integrando se anula donde $z=0$.
En lo físico cada medida basal asigna masa $1/j$ a cada índice.
La variación del integrando en $v$ es
$-\log(v/z)-2-\log(1-\sigma t_rt_c v)$, una suma de potenciales por la ecuación
de equilibrio. Así el término de la variación del óptimo se anula al derivar en
$\sigma$ con márgenes fijos. Esta es una identidad entre las prolongaciones:
las variaciones satisfacen margen cero en cada lado aun para las medidas complejas.
Con $E_{ab}=\int z_x^2\,dM_a^R dM_b^C$, se deduce en los parámetros físicos
\[
 R(x)={\rm af\acute in}(x)+s^{-1}\mathcal R_0(sx)+\tfrac12 E_{22}
                                  +s\,\Psi_R(s,h,x).
\]
La notación $sx$ significa las cuatro perturbaciones con signo por sus respectivos átomos.
$\mathcal R_0$ es el valor a $\sigma=0$: es analítico y acotado en esas perturbaciones.
El término lineal en $\sigma$ es $E_{22}/2$ antes de dividir por $s$.
$\Psi_R$ es analítico y acotado por Taylor en $\sigma$ en la familia uniforme del apartado 3.
Cuando $h_g=0$ cualquier integral pierde los índices especiales del grupo $g$.
Las ecuaciones restringidas a los demás índices y sus integrales son independientes de
los $x_i$ de ese grupo. Esta propiedad se conserva en cada resto de Taylor.
En consecuencia, la parte de $\Psi_R$ dependiente de ambos grupos después de centrar en
cada grupo es $O(h_1h_2)$ en todo el polidisco. Se obtiene por dos aplicaciones de Cauchy
o de la fórmula integral de diferencias.
